\begin{frame}{Agents and LLMs: Covered in the Advanced NLP Track}

    Everything in this lecture applies to a model that maps an input to a prediction. Systems that
    \textbf{take actions} --- call tools, browse, write files, spend money --- add a distinct threat
    model that this deck deliberately does not duplicate.

    \vspace{0.9em}

    \begin{columns}[T]
        \begin{column}{0.49\textwidth}
            \begin{tcolorbox}[colback=gray!8, colframe=gray!60,
                              title={Phase 2B: AI Safety for Agents}]
                \scriptsize
                \begin{itemize}
                    \item Prompt injection, direct and indirect.
                    \item Excessive agency and tool misuse.
                    \item Sandboxing, least privilege, and human-in-the-loop gates.
                    \item Agent-specific benchmarks and red-teaming.
                    \item OWASP Top 10 for LLM Applications.
                \end{itemize}
            \end{tcolorbox}
        \end{column}
        \begin{column}{0.49\textwidth}
            \begin{tcolorbox}[colback=raiBlue!5, colframe=raiBlue!60, title=What carries over]
                \scriptsize
                \begin{itemize}
                    \item Disaggregated evaluation --- an agent can fail more often for some users
                          or some languages.
                    \item Documentation of intended and out-of-scope use.
                    \item The same habits either way: a named reviewer, a stated scope, incident
                          reporting, and a re-audit date apply to an agent exactly as they do to a
                          classifier.
                \end{itemize}
            \end{tcolorbox}
        \end{column}
    \end{columns}

    \vspace{0.7em}
    \centering
    \textit{If you are shipping an agent, take both lectures. Neither is a substitute for the other.}
\end{frame}
